\begin{ExerciseBlock}
\ExerciseHeading{03}{Scheinbare Lage einer Münze im Wasser}{geometrical_and_wave_optics_apparent_coin_position}
Sie beobachten eine Münze auf dem ebenen Boden eines Wasserbeckens der Tiefe $h=\qty{1}{\metre}$. Der beobachtete Lichtstrahl bildet in Luft mit dem Lot den Winkel $\alpha=\ang{45}$. Die Brechungsindizes sind $n_\mathrm W\approx1{,}33$ und $n_\mathrm L\approx1$.
\begin{center}
\CachedTikzFigure{content/resources/chapters/geometrical_and_wave_optics/tikz/coin_in_water.tex}
\end{center}
Berechnen Sie den horizontalen Abstand $x$ zwischen der Münze und der Stelle am Beckenboden, an der die rückwärtige Verlängerung des beobachteten Strahls den Boden schneidet.
\end{ExerciseBlock}
\UseExerciseSolution{geometrical_and_wave_optics_apparent_coin_position}
